\section{Current dependency audit: original global consumers, theta derivatives, and Count-II normalization}
\label{sec:dependency-audit-v34}

\status{Version 0.34 accepts the underlying Alp\"oge--Claude formal
proof as fully checked and correct. It adds 55 human proof claims:
GVC1--18, GFA1--12, TAE1--13 and CIP1--12. The human reconstruction
remains incomplete. Completed implications retain their original
producer premises; category D concerns our understanding of the
formal source. This current audit supersedes historical frontier
statements without changing any historical mathematical proof.}

\subsection{What this continuation supplies}

\paragraph{The exact global numerical consumers.}
GVC reconstructs the degree products and their finite Vandermonde
comparison on the original words. The large top--bottom factors
cancel in the cross-word ratio because the actual multiplicities
and shifted coordinates match. The remaining factors are controlled
uniformly by the original parameter windows. Together with the
previous mass and index identities, this proves the literal
\texttt{CountF}. It also installs all sixteen numerical and eight
weight clauses of \texttt{NumsL}/\texttt{WtsL}, with one threshold
before the word and family index.

Direct sum estimates applied to the original lower and maximal
column slopes give the exact \texttt{Conclusion5}/\texttt{Conclusion6}
and \texttt{F5}/\texttt{F6} consumers, and then
\texttt{SlopesMF}. Every original determinant shift is retained.
These sufficient arguments do not automatically prove every stronger
intermediate statement called \texttt{PartFive} or \texttt{PropZ}
in the source. GVC18 also checks that the needed bound on
\(L.\mathrm{logDerivativeAtOne}\) is the \texttt{DlSize} projection
of the original \texttt{FamP.sizes}; it is not a new independent
producer for these consumers. Existence of that original family
remains a separate obligation.

\paragraph{Quotient slopes with the actual type-prime cost.}
For the original lattice \(N_u=I_uN_u^0\), GFA1 proves the precise
congruence cost
\[
 \frac1{r_u}\sum_{w\mid T_u}2\ell_we_w
                  L.\mathrm{lenMN}(u,w)
 \le36|TT|m(a/k+b')\log T_u.                         \tag{AU34.1}
\]
A finite subset of the places is allowed. The weight calculation
first gives constant nine, and the original retained-rank ratio is
at least one half. The source's factor two in the squared metric
then gives exactly 36. The original local length and ramified rank
identities, not substitute lattices, enter this calculation.

For the correction away from that prime, a useful direct argument
uses only what the consumer requires. Put
\[
 b=e^{x^\epsilon},\quad a=xb,\quad
 X'=e^{x^{\epsilon/3}},\quad S=e^{x^{\epsilon/2}}.
\]
The actual definitions of the local type-I/type-II correction and
ceilings give \(\tau_p\le aS^{20}\), with \(\tau_p=0\) for
\(p\ge X'\). The finite-place degree identity and a crude count of
primes then give
\[
 T_a=\sum_w2\ell_we_w\tau_{p(w)}
       \le aS^{24}\le b^2                              \tag{AU34.2}
\]
for all sufficiently large \(x\), uniformly before every index.
This bypasses the stronger weighted-prime estimate solely for this
consumer. It does not assert that stronger estimate has been proved.
The exact ideal shift, integral local quotient lengths, nonnegative
rounding error and the zero contribution at the exceptional prime
now give \texttt{PlacesF} and \texttt{SlopesNF} in the source's
normalization, from the original local-gain and metric premises.

\paragraph{The frame inequality has its full height quantifier.}
GFA6--9 check positive target rank, the algebraic dual with its
original metric, the actual rational Bessel map and every Gram
normalization. Let \(s\) be the weighted mean source slope minus
height, let \(u\ge0\) be the maximal excess above \(s\), let
\(\delta\) be the target discard, and let \(d,e_0\) be the
original source-discard and diagonal errors. The reconstructed
frame estimate gives, for \emph{every} \(H>0\) with
\(s-\mu_{\min}(RR^\vee)\le H\),
\[
 \mu(RR^\vee)\ge s-
 \frac{4uH/(u+H)+(5\sigma+4d+4\sqrt{e_0})H}{1-\delta}.
                                                               \tag{AU34.3}
\]
For \(u>0\), \(H/(u+H)\le\log(1+H/u)\); for \(u=0\), both
terms multiplied by \(u\) are zero. This proves the literal
\texttt{Frame\_at} with its logarithmic term. No upper restriction
on \(H\) is added. One pair of error bounds is chosen before the
family index and yields \texttt{FrameF} with the original \(C_H\).

GFA10--12 check the actual node codes and intervals, all seventeen
\texttt{Win\_at} clauses, and all four \texttt{Char\_at} clauses.
The primitive-character choices agree at all positive integers,
including ramified ones. Uniqueness uses CRT, multiplicativity and
the minimum conductor at each prime; it needs no theorem asserting
a prime in each arithmetic progression. The three quadratic
readings use the original host's stored \texttt{flPS} field, with
the correct interchange of the two sign coordinates. Thin-character
and quadratic estimates remain the original chosen inputs.

The bounded residual global task is now explicit: prove the
original \(\texttt{LossFS}(\mathcal F,0)\), and then carry the
same selected values and characters through the actual
\texttt{linkCore}/\texttt{ChainTied} and final \texttt{Closed}
wrapper. The constants remain
\(C_{\rm abs}=2,C_0=3,C_5=4,e_D=2,e_L=0\).
The compatible second-layer selector and local/spectral premises
are not supplied by these consumer proofs.

\paragraph{Real derivatives and one growth exponent for the original theta integral.}
TAE first differentiates the reconstructed Gaussian oscillator
formula in every Heisenberg smooth seminorm. Polynomial Gaussian
domination justifies the derivatives at every order. The actual
full action agrees with that oscillator construction on the
original theta space, with its continuous scalar character.
For an arbitrary adelic left factor, exact Haar \(L^2\) unitarity
transports the derivative; continuity converts equality almost
everywhere to equality at every point. No unproved supremum-norm
boundedness of that left factor is assumed.

A common finite Heisenberg level thickens the real cube into an
adelic set of positive measure. Compactness bounds its rational
translate multiplicity, and unitarity bounds its integral.
Iterating a one-dimensional evaluation inequality then proves
\[
 |R_{\mathbf X}\Phi_L(xg)|\le C_{\mathbf X,x}
                       H_T(g)^N,
 \qquad N=\dim_{\mathbb R}\mathrm{RW}(T,S),             \tag{AU34.4}
\]
with the \emph{same} \(N\) for every derivative word. The constants
may depend on the word and translate. The source has the original
compact rational representatives from BTH3; a finite coefficient
basis supplies a dominating function and justifies differentiation
under the same integral and Haar measure.

TAE also reconstructs the stored-vacuum passage, using the original
universal finite place-model condition. Adding an algebraically
extended annihilator to the place-model map leaves its defining
finite intertwiners and initial vacuum unchanged. Universality
therefore forces the stored vector to satisfy the real vacuum
equations. The exact commutator computation then identifies the
complex lowering generator with the original mixed Laplacian,
including coefficient \(-4\pi i/s_\tau^2\). Positive and negative
minors in the literal \(\Delta\) use disjoint source columns, so
that Laplacian kills \(\Delta\).

Under the explicitly retained lowest-type smoothness and
irreducibility premises, the source and compact target actions
propagate lowering from \(\Delta\) to every coefficient value.
Combined with FIA/CTC cuspidality and the original finite-level
identities, this gives all eight fields of
\texttt{IsHolCuspKMap} for the exact lowest-type integral.
It is not yet the required Hilbert multiplicity map.
The remaining C subcontracts are: one \(io_i\) valid for every
parameter \(t\); \texttt{LowSm}/\texttt{LowIrr} for the full
original lowest compact type, including multiplicity one;
boundedness and target finite volume; and the original extremal
coefficient/Hilbert realization with Sobolev continuity. Same-Bundle
readings, rational tests and actual base rows remain separate.

\paragraph{Exact reciprocity and the required Count-II normalization.}
CIP matches classical local/global reciprocity to every source
predicate map in the required split-CM scope. Arithmetic reciprocity
sends a uniformizer to arithmetic Frobenius, whereas the source
requires the inverse valuation sign. The proof explicitly inverts
that convention. Global-map uniqueness is checked in every finite
quotient. At a split CM place \(w\ne c(w)\), a local unit in the
global reciprocity kernel is in the finite-idelic closure of the
global units. For any global unit \(e\), \(e/c(e)\) is a root of
unity; its finite possibilities and the coordinates away from
\(w,c(w)\) force the original local unit to be one. This proves
the required equality for \emph{every} \texttt{IsLocalArtin},
not only a constructed witness. Norm compatibility and the checked
inversion give the source's positive cyclotomic norm sign.
The three category-A application certificates cite Milne,
\emph{Class Field Theory}, version 4.03, with exact theorem locators
and checked number-field/local-field hypotheses. The convention,
uniqueness and character matching are proved separately in CIP.

For the irreducible-point normalization, put
\(A=K[[t]]\) in odd characteristic. An integral unimodular symmetric
form has a vector of unit self-pairing: otherwise polarization,
using the invertibility of two, would contradict unimodularity.
Integral Gram--Schmidt diagonalizes it over \(A\). A finite extension
of the constant field supplies the finitely many residue square
roots, and coefficient recursion supplies square roots of the
resulting unit power series. This gives \(P=1\) by an integral
basis change while keeping \(t\) unchanged.

The transformed representation remains integral. Absolute
irreducibility gives full matrix span; intersections with the
original invariant flag give saturated integral submodules in the
same order. The proof checks the original characters, polynomial
law, determinant, N3 and N4, including the \emph{old} \(D_0\).
Thus the entire universal \texttt{NiceI} instance of \texttt{Norm}
required by JSP11 is supplied, including characteristic three.
The general \texttt{NiceE} normalization is not asserted.
The ordinary arithmetic/control system, four further Hida fields,
same-ideal all-level \texttt{HasAut}, and the other single-point
support inputs remain C producers. Both original rank-four Links
still use one coefficient threshold at fixed \(\iota\).

\subsection{The selected complete-proof cut and next priorities}

The final interface is unchanged: one compatible design and nine
original \texttt{Closed} predicates. The cutoff and transport tuple
are joined before choosing the common family; eventual thresholds
come before its indices, words and places. RH and two-primary
targets require I/II/V; unrestricted Dirichlet GRH additionally
requires III, SeedIV and the original chains. The two fixed fields
in Counts I and V remain distinct. Keep the exact host window
\(\mathrm{VWinAt}(\chi,401)\), cutoff \(\lfloor x^{400}\rfloor\),
and final numerical exponent \(1/2+9\epsilon\). The AQ alternative
retains its own arithmetic envelope.

The selected graph still has 35 open coupled leaves: 23 C, 10 B
and 2 D. This unchanged count is compatible with substantial
progress inside the three refined leaves. The nested Count-II
reciprocity-understanding issue is now closed; it was not an
additional top-level D leaf. The 121-node graph, 245 historical
targets in 107 groups, thirteen shared refinements and eight
optional routes remain available. None of these counts is a
percentage of proof completion or an absolute minimality claim.

The next bounded global work is the original loss estimate, then
its exact tied-ledger and Closed installation. In parallel, the
common \(ioP\), lowest-type multiplicity-one and Hilbert realization
serve several analytic branches; the ordinary rank-four
arithmetic/control and same-ideal automorphic producer remains an
independent deep obstruction already needed for RH. The recursive
common construction, local \texttt{NBuilt}/\texttt{NdSpec}, and the
original denominator, retained defect, period-height, local-gain
and spectral producers must stay visible as shared dependencies.

A means a direct published application with hypotheses checked;
B an elementary or moderately difficult independent reconstruction;
C a substantial deep reconstruction; D a source-understanding
issue in the human reconstruction. Completed A applications remain
in the registry. The following register gives every remaining
selected contract verbatim, without replacing common witnesses
by unrelated existential choices.
